Nền tảng kỹ thuật thứ ba của CausClass là sử dụng LLM để hỗ trợ tìm kiếm và diễn giải đồ thị. LLM có khả năng xử lý mô tả biến bằng ngôn ngữ tự nhiên và khai thác tri thức miền ở mức khái quát. Trong bài toán lớp học, tri thức này có thể giúp nêu các giả thuyết như hành vi trao đổi có thể liên quan đến các hành vi tương tác khác, hoặc hành vi sử dụng điện thoại có thể tạo ra mẫu thời gian khác với đọc và viết. Tuy nhiên, các giả thuyết như vậy chỉ có giá trị nếu được kiểm chứng bằng dữ liệu quan sát.

Các nghiên cứu gần đây thường dùng LLM theo hai cách. Cách thứ nhất xem LLM như nguồn prior cấu trúc: mô hình gợi ý hướng cạnh, quan hệ hợp lý hoặc ràng buộc giữa biến, sau đó thuật toán thống kê dùng các gợi ý này để thu hẹp không gian tìm kiếm \,\cite{darvariu2024llmpriors,barakati2025llmassisted}. Cách thứ hai dùng LLM như tác nhân thăm dò đồ thị, trong đó mô hình đề xuất các bước chỉnh sửa dựa trên đồ thị hiện tại, phản hồi từ dữ liệu và mô tả nhiệm vụ \,\cite{jiralerspong2024efficient,roy2025causalllm,peng2026causcientist}. Cả hai cách đều cho thấy LLM có thể hỗ trợ tìm kiếm, nhưng không loại bỏ nhu cầu đánh giá định lượng.

Trong CausClass, LLM được đặt vào vai trò bộ sinh đề xuất. Ở mỗi vòng tìm kiếm, LLM nhận mô tả biến hành vi, đồ thị hiện tại, danh sách tabu, phản hồi từ các chỉnh sửa bị loại và tín hiệu từ AERCA, chẳng hạn các cạnh mạnh nhưng chưa được giữ trong đồ thị hiện tại. Mô hình có thể đề xuất thêm hoặc xóa một số cạnh. Sau đó, mỗi chỉnh sửa được chuyển thành đồ thị ứng viên, áp mask vào AERCA và đánh giá lại bằng lỗi kiểm chứng cùng phạt độ phức tạp. Như vậy, LLM mở rộng không gian giả thuyết, còn verifier định lượng quyết định ứng viên nào được giữ.

Chiến lược tìm kiếm được dùng là beam search. Thay vì chỉ giữ một đồ thị tốt nhất tại mỗi bước như tìm kiếm tham lam, beam search giữ một tập nhỏ các đồ thị hứa hẹn để tiếp tục mở rộng. Điều này giảm nguy cơ mắc kẹt ở một chỉnh sửa cục bộ. So với tìm kiếm ngẫu nhiên, LLM giúp ưu tiên các chỉnh sửa có ý nghĩa miền trước khi đưa vào kiểm chứng. Cơ chế này phản ánh quan điểm thiết kế của đồ án: tri thức ngữ nghĩa giúp định hướng, còn dữ liệu quyết định.

LLM cũng được dùng trong phần sinh báo cáo quan sát. Báo cáo cần chuyển thống kê chuỗi thời gian và danh sách cạnh thành văn bản dễ đọc, nhưng không được đổi chiều cạnh, phát minh cạnh mới hoặc biến liên kết dự báo thành kết luận nhân quả. Vì vậy, CausClass render xác định các phần dữ kiện như thống kê hành vi và danh sách cạnh trước, sau đó chỉ cho LLM viết các đoạn diễn giải có ràng buộc. Báo cáo phải bắt đầu từ bằng chứng phát hiện, trình bày liên kết thời gian như giả thuyết cần kiểm tra và khuyến nghị người dùng quay lại video khi cần đưa ra nhận định.

Tóm lại, Chương 2 đã trình bày ngữ cảnh bài toán, các nghiên cứu tương tự và những nền tảng kỹ thuật chính của CausClass. Các nội dung này dẫn đến thiết kế ở Chương 3: YOLO26n-MHSA đảm nhiệm perception thị giác, ByteTrack và time bin tạo chuỗi hành vi, AERCA kiểm chứng đồ thị liên kết dự báo, LLM đề xuất chỉnh sửa cấu trúc và bộ sinh báo cáo trình bày kết quả theo nguyên tắc evidence-first. Cách tổ chức này giúp hệ thống cân bằng ba yêu cầu: tự động hóa, khả năng diễn giải và sự thận trọng khi xử lý dữ liệu giáo dục quan sát.